\subsection{Generalities}

Recall (cf.~TG, IX, p.~37, def. 9) that a normed algebra is an algebra A over K equipped with a norm \(x \mapsto \|x\|\) such that:

\[
\| x y \| \leqslant \| x \| \| y \|\tag{1}
\]

for all \(x, y \in A\) . If A is complete, A is said to be a Banach algebra.

Let A be a complete normable algebra over K. The topology of A can be defined by a norm satisfying (1) (cf.~TG, IX, p.~38). When A is equipped with the structure of a normed algebra defined by such a norm, we shall also say that the algebra A is a Banach algebra.

We also recall the following facts (cf.~TG, IX, p.~38--39):

\begin{enumerate}
\def\labelenumi{(\arabic{enumi})}
\item
  If A is a nonzero normed algebra and has a unit element e, then \(\|e\|\geqslant1\) .
\item
  Let A be a normed K-algebra. The opposite algebra of A, equipped with the same norm, is a normed algebra. The completion of A, equipped with the norm obtained by continuous extension of the norm of A, is a K-Banach algebra. Every subalgebra of A, equipped with the induced norm, is a normed algebra. If I is a closed two-sided ideal of A, the algebra A/I, equipped with the norm defined by \(\|\dot{x}\| = \inf_{x \in \dot{x}} \|x\|\) for all \(\dot{x} \in A/I\) is a normed algebra.
\item
  Let \((\mathrm{A}_{i})_{i\in\mathrm{I}}\) be a family of normed algebras and B the product algebra of the algebras \(A_{i}\) . The subalgebra of elements \((x_{i})_{i\in\mathrm{I}}\) of B such that \(\|(\boldsymbol{x}_{i})\|=\sup_{i\in\mathrm{I}}\|\boldsymbol{x}_{i}\|<+\infty\) is a normed algebra, called the product normed algebra of the algebras \((\mathrm{A}_{i})_{i\in\mathrm{I}}\) . If \(A_{i}\) is a Banach algebra for every i, then B is a Banach algebra.
\end{enumerate}

Let A be a normed algebra. Let \((y_{i})_{i\in\mathrm{I}}\) be a family of elements of A; the smallest closed subalgebra B of A containing the elements \(y_{i}\) is called the closed subalgebra of A generated by the \(y_{i}\) ; if B = A, one says that the \(y_{i}\) topologically generate the normed algebra A, or that the family \((y_{i})_{i\in\mathrm{I}}\) is a topological generating system of the normed algebra A.

Similarly, if A is a unital normed algebra, the smallest closed unital subalgebra containing the elements \(y_{i}\) is called the closed unital subalgebra of A generated by the \(y_{i}\) . If it is equal to A, one says that the \(y_{i}\) topologically generate the unital normed algebra A.

Let A be a normed K-algebra. Denote by \(\tilde{\mathbf{A}}\) the algebra obtained from A by adjoining an identity element. On \(\tilde{\mathbf{A}}\) , one defines a norm by setting \(\| (\lambda ,x)\| = |\lambda | + \| x\|\) for all \(\lambda \in \mathbf{K}\) and every \(x\in \mathbf{A}\) . We have

\[
\begin{array}{r l} & {\| (\lambda , x) (\mu , y) \| = | \lambda \mu | + \| x y + \mu x + \lambda y \|} \\ & {\qquad \leqslant | \lambda | | \mu | + \| x \| \| y \| + | \mu | \| x \| + | \lambda | \| y \|} \\ & {\qquad = \| (\lambda , x) \| \| (\mu , y) \|.} \end{array}
\]

Thus \(\widetilde{\mathsf{A}}\) becomes a normed algebra, called the unital normed algebra deduced from A by adjoining an identity element. The algebra A is identified with the closed two-sided ideal \(\{0\} \times \mathbf{A}\) of \(\widetilde{\mathsf{A}}\) .

DEFINITION 1. --- Let A be a normed algebra. For \(a \in \mathrm{A}\) , let \(\gamma_{a}\) and \(\delta_{a}\) be the maps \(x \mapsto ax\) , and \(x \mapsto xa\) from A to A. The map \(\gamma: a \mapsto \gamma_{a}\) is a representation of A in A, called the left regular representation of A. The map \(\delta: a \mapsto \delta_{a}\) is a representation of the opposite algebra of A in A, called the right regular representation of A.

Lemma 1. --- Let A be a normed algebra. The left regular representation of A and the right regular representation of A are continuous with norm \(\leqslant 1\) .

If the algebra A is unital, the maps \(x \mapsto \|\gamma_{x}\|\) and \(x \mapsto \|\delta_{x}\|\) are norms on A, defining the topology of A. They satisfy the inequality (1).

If the unital normed algebra A is nonzero, that is, if \(e \neq 0\) , we also have \(\|\gamma_{e}\| = \|\delta_{e}\| = 1\) .

It is immediate that \(\|\gamma_{x}\|\leqslant\|x\|\) and that \(\|\delta_{x}\|\leqslant\|x\|\) .

If A has an identity element \(e\) , then we have \(x = \gamma_x e = \delta_x e\) , hence:

\[
\| x \| \leqslant \| \boldsymbol {\gamma} _ {x} \| \cdot \| e \| \quad \| x \| \leqslant \| \boldsymbol {\delta} _ {x} \| \cdot \| e \|.\tag{2}
\]

In this case, \(x \mapsto \|\gamma_{x}\|\) and \(x \mapsto \|\delta_{x}\|\) are therefore norms equivalent to the norm of A. They satisfy (1) because \(\gamma\) and \(\delta\) are representations.

The last assertion follows from the fact that \(\gamma_{e} = \delta_{e} = Id_{A}\) .

